\subsection{逆関数の微分}
一つの区間で連続かつ狭義単調な関数 $y=f(x)$ は,その値域で逆関数 $x=g(y)$ をもつ.このとき逆関数も連続であるという基本性質を用いる.その区間の内点 $a$ について $b=f(a)$,\,$f'(a)\ne0$ とし,\,$y=b+k$ に対応する $x$ を $a+h$ と書けば,\,$k\to0$ のとき $h\to0$ であり,\,$k=f(a+h)-f(a)$ である.したがって
\[
 \frac{g(b+k)-g(b)}k
 =\frac{h}{f(a+h)-f(a)}
 =\cfrac{1}{\cfrac{f(a+h)-f(a)}h}
 \longrightarrow\frac1{f'(a)}.
\]
よって
\[
 g'(b)=\frac1{f'(a)},\qquad
 g'(y)=\frac1{f'(g(y))}
\]
となる.分母が $0$ でないことが必要である.たとえば $f(x)=x^2$ を $x>0$ に限ると,逆関数は $g(y)=\sqrt y$ であり,\,$g'(y)=\frac1{2\sqrt y}$ となる.\,$x^2$ を実数全体で考えると,正負の二つの $x$ が同じ値に対応して逆関数が定まらないので,先にどちらの枝を使うかを決めなければならない.


